\documentclass[11pt,a4paper]{altacv}

\geometry{left=1.5cm,right=1.5cm,top=1.cm,bottom=1.2cm,columnsep=1.5cm}

\usepackage[utf8]{inputenc}
\usepackage[T1]{fontenc}
\usepackage[french]{babel}
\usepackage{paracol}
\usepackage{xcolor}
\usepackage{hyperref}
\usepackage{fontawesome5}
\usepackage{setspace}

\definecolor{SlateGrey}{HTML}{2E2E2E}
\definecolor{LightGrey}{HTML}{666666}
\definecolor{DarkPastelRed}{HTML}{8AC0D9}
\definecolor{PastelRed}{HTML}{00B0DA}
\definecolor{GoldenEarth}{HTML}{B4AD0C}
\colorlet{name}{black}
\colorlet{tagline}{PastelRed}
\colorlet{heading}{DarkPastelRed}
\colorlet{headingrule}{GoldenEarth}
\colorlet{subheading}{PastelRed}
\colorlet{accent}{PastelRed}
\colorlet{emphasis}{SlateGrey}
\colorlet{body}{LightGrey}

\renewcommand{\familydefault}{\sfdefault}

\name{\Large Guillaume BOILEAU}
\tagline{PMO Fundings / CIR | R\&D, financements innovation, coordination internationale \& rédaction technique}
\personalinfo{
  \email{guillaume.boileaupro@gmail.com}
  \phone{+33 6 59 94 85 84}
  \location{Alpes-Maritimes, France}
  \linkedin{guillaume-boileau-phd-891b81265}
}

\setlength{\parskip}{0.75\baselineskip}

\begin{document}
\makecvheader
\vspace{-0.6cm}

\cvsection{Lettre de motivation}
\vspace{-0.4cm}

{\color{accent}\textbf{Objet :}} \textit{Candidature -- PMO Fundings -- Sophia Antipolis}

{\color{body}

Madame, Monsieur,

Je vous adresse ma candidature pour le poste de \textbf{PMO Fundings}. Votre offre m'intéresse particulièrement car elle combine plusieurs dimensions directement cohérentes avec mon parcours : coordination transverse, environnement international, projets d'innovation, compréhension de sujets Cloud/IA/Data/cybersécurité, rédaction de documents techniques exigeants, Crédit Impôt Recherche et financements publics.

Docteur en physique, j'ai construit mon parcours dans des environnements R\&D à forte complexité technique, où il faut comprendre rapidement les enjeux scientifiques, identifier les verrous, structurer les hypothèses, consolider les contributions, documenter les résultats et produire des synthèses exploitables par des interlocuteurs variés. Cette capacité à faire le lien entre expertise technique, documentation rigoureuse et coordination transverse correspond au cœur du poste proposé.

Au \textbf{CNES}, j'ai travaillé sur le segment sol de la mission spatiale \textbf{LISA}, dans un environnement CNES/ESA associant données mission, niveaux L0--L3, cycle en V, simulation, intégration de modèles et validation technique. J'ai notamment développé \textbf{CATpyLISA}, une plateforme Python destinée à structurer des workflows de données, comparer des modèles, contrôler la cohérence des résultats et préparer des éléments de revue. Ce travail m'a amené à dialoguer avec plusieurs contributeurs, à formaliser des méthodes, à analyser des écarts, à documenter des limites et à produire des synthèses techniques précises.

Cette expérience est directement transposable aux activités CIR et financements innovation : conduite d'interviews techniques, collecte d'informations auprès d'équipes projet, identification de verrous scientifiques ou technologiques, formalisation de l'état de l'art, description de la démarche suivie, consolidation de résultats et rédaction de notices techniques. Même si mon parcours n'est pas initialement celui d'un consultant CIR, je dispose d'une pratique concrète de la R\&D, de la rédaction scientifique et de la structuration de dossiers techniques à haut niveau d'exigence.

J'ai également travaillé sur des sujets combinant \textbf{IA, data science, logiciel, simulation et validation}. À l'Universiteit Antwerpen, j'ai développé des workflows Python pour analyser des données capteurs, modéliser des signaux faibles, comparer des configurations et produire des recommandations techniques. Plus récemment, chez \textbf{VEV Platform Services France}, j'ai renforcé ma capacité à comprendre un environnement logiciel et data opérationnel, analyser des incidents, collecter des informations techniques et restituer clairement des constats exploitables.

Pour Scalian, je peux apporter un profil hybride \textbf{R\&D, data/IA, logiciel, coordination et rédaction technique}. Je suis capable de comprendre des sujets complexes, de dialoguer avec des équipes techniques, de structurer des informations dispersées, de suivre plusieurs actions en parallèle, de produire une documentation claire et de travailler dans un contexte international. Je peux également contribuer à l'harmonisation des pratiques, à la consolidation de contributions issues de différentes Business Units et à la préparation de dossiers de financements européens ou internationaux.

Je souhaite aujourd'hui mettre cette expérience au service d'un rôle PMO Fundings, à l'interface entre innovation, financement, conformité, coordination et valorisation technique. Le poste proposé par Scalian représente pour moi une continuité logique vers des activités où la rigueur scientifique, la capacité de synthèse, la compréhension technologique et le pilotage transverse sont déterminants.

Je serais heureux de pouvoir échanger avec vous afin de vous présenter plus en détail mon parcours et la manière dont mon expérience R\&D, IA/Data, documentation technique et coordination internationale peut contribuer efficacement à vos activités de financements publics, CIR et dispositifs d'innovation.

Je vous prie d'agréer, Madame, Monsieur, l'expression de mes salutations distinguées.

\textbf{Guillaume Boileau}

}

\end{document}
